\section{Useful lemmas}

%%%%%%%%%%%%%%%%%%%%%%%%%%%%%%%%%%%%%%%%%%%%%%%%%%%%%%%%%%%%%%
\begin{lemma}\label{lem:svd}
Let $B^k$ be the unit $k$-ball in $\R^k$, let
$L:\mathbb{R}^k\longrightarrow\mathbb{R}^m$ be a linear map of rank
$m\le k$.  Then $\HH^m(L(B^k)) = \omega_m \sqrt{\det(LL^T)}$.
\end{lemma}

\begin{proof}
Let $L=U\Sigma V^T$ be a singular value decomposition of $L$, where
$U\in O(m)$, $V\in O(k)$, and
$\Sigma=\operatorname{diag}(\sigma_1,\ldots,\sigma_m,0,\ldots,0)$.

Since orthogonal transformations preserve Euclidean balls, we have
$V^T B^k = B^k$. Since they also preserves $m$-dimensional volume, we
have $\HH^m(L(B^k)) =
\HH^m(\Sigma(B^k))$.

%% The image $\Sigma(B^k)$ is the ellipsoid with semi-axes
%% $\sigma_1,\ldots,\sigma_m$, whose volume is $\omega_m \prod_{i=1}^m
%% \sigma_i$.
\begin{sloppypar}
Now $\Sigma(B^k)= \left\{ (\sigma_1x_1,\ldots,\sigma_mx_m):
x_1^2+\cdots+x_k^2\le1 \right\}$.  Its image is precisely the
ellipsoid $\left\{ y\in\mathbb{R}^m:
\sum_{i=1}^m\frac{y_i^2}{\sigma_i^2}\le1 \right\}$, whose semi-axes
are $\sigma_1,\ldots,\sigma_m$. Its volume is $\omega_m \prod_{i=1}^m
\sigma_i$.
\end{sloppypar}

On the other hand, $L L^T = U \Sigma \Sigma^T U^T$. Hence $\det(L
L^T)=\prod_{i=1}^m \sigma_i^2$.  Gathering all identities gives

\begin{minipage}[b]{0.9\textwidth}
\begin{align*}
  \HH^m(L(B^k)) = \HH^m(\Sigma(B^k))
  = \omega_m \prod_{i=1}^m \sigma_i = \omega_m \sqrt{\det(LL^T)}.
\end{align*}
\end{minipage}
\qed
\end{proof}

\section{Geometric invariants for balls and half-balls}

Here is an illustration of the geometric invariants $\bar{W}_{m}(k)$
of the $k$-ball for different $k$ and $m$.  We also displays the
values $\bar{W}^{1/2}_{m}(k)$ corresponding to the geometric
invariants for the half ball of dimension $k$.
\begin{center}
\begin{tabular}{|c|ccccc|} \cline{2-6}
  \multicolumn{1}{c}{$\bar{W}_{m}(k)$} & \multicolumn{5}{|c|}{$m$} \\ \cline{1-1}
  $k$ & 0 & 1 & 2 & 3 & 4 \\ \cline{2-6}
  0   & 1 &   &   &   &   \\
  1   & 1 & 4 &   &   &   \\
  2   & 1 & 8 & $\pi^2$  &   &   \\
  3   & 1 & 12 & $3\pi^2$  & $\frac{16}{9}\pi^2$   &   \\
  4   & 1 & 16 & $6\pi^2$  & $\frac{64}{9}\pi^2$   & $\frac{1}{4}\pi^2$  \\ \hline
\end{tabular}
\quad
\begin{tabular}{|c|ccccc|} \cline{2-6}
  \multicolumn{1}{c}{$\bar{W}^{1/2}_{m}(k)$} & \multicolumn{5}{|c|}{$m$} \\ \cline{1-1}
  $k$ & 0 & 1 & 2 & 3 & 4 \\ \cline{2-6}
  0   & 1 &   &   &   &   \\
  1   & 1 & 1 &   &   &   \\
  2   & 1 & 5 & $\frac{1}{4}\pi^2$  &   &   \\
  3   & 1 & 9 & $\frac{3}{2}\pi^2$  & $\frac{4}{9}\pi^2$   &   \\
  4   & 1 & 13 & $\frac{15}{4}\pi^2$  & $\frac{28}{9}\pi^2$   & $\frac{1}{16}\pi^2$  \\ \hline
\end{tabular}
\end{center}
